\section{Мета роботи та завдання}

\subsection{Мета лабораторної роботи}

Метою лабораторної роботи є отримати практичні навички зі створення і використання 
абстрактних класів та інтерфейсів, засвоїти принципи успадкування класів у мові 
програмування Java, навчитися створювати ієрархії класів, перевизначати методи 
базових класів, реалізовувати методи інтерфейсів та використовувати поліморфізм 
під час роботи з об’єктами. У процесі виконання роботи необхідно:

\subsection{Загальне завдання (Практична частина)}

Під час виконання лабораторної роботи необхідно створити Java-програму, у якій 
буде продемонстровано використання успадкування, абстрактного класу, інтерфейсу 
та поліморфізму.:
\begin{enumerate}
    \item Створення базового абстрактного класу;
    \item Створення дочірніх класів;
    \item Створення інтерфейсу;
    \item Створення об’єктів;
    \item Демонстрація поліморфізму;
    \item Робота з інтерфейсом;
    \item Розширення програми:
    Після виконання базової частини необхідно адаптувати приклад до індивідуального завдання.
Програма повинна містити:
    \begin{enumerate}
        \item щонайменше один абстрактний клас;
        \item щонайменше два дочірні класи;
        \item щонайменше один інтерфейс;
        \item реалізацію інтерфейсу конкретним класом;
        \item перевизначення методів;
        \item конструктори базового та дочірніх класів;
        \item використання \texttt{super};
        \item демонстрацію поліморфізму;
        \item створення та використання об’єктів у методі \texttt{main()}.
    \end{enumerate}

\end{enumerate}

\subsection{Індивідуальне завдання. Варіант №7. Комп'ютерна техніка.}

\begin{enumerate}
    \item Створити абстрактний клас \texttt{ComputerDevice} з полями \texttt{brand} та \texttt{model};
    \item Створити класи \texttt{Laptop} та \texttt{Desktop};
    \item Додати інтерфейс \texttt{Connectable} з методом \texttt{connect()}.
\end{enumerate}

\subsection{Оформити звіт і відповісти на контрольні питання.}


